\appendix
\section*{Appendix R5: Interpretive Field Bundle Formalization}

\subsection*{R5.0 Overview}

This appendix provides a formal construction of the interpretive-field
bundle $E$ used in the NDH 42D Interpretation Algebra. The bundle
structure ensures that interpretive-field operators are mathematically
well-defined as smooth sections, and that the interpretation operator
$I_p$ interacts consistently with the geometric and algebraic structure
of the 42D interpretation manifold $X_{42}$.

\subsection*{R5.1 Base Manifold}

Let $X_{42}$ be the 42-dimensional smooth manifold constructed in
Appendix R4, with maximal smooth atlas $\mathcal{A}$.

\subsection*{R5.2 Fiber Structure}

Let $F$ be a smooth manifold representing interpretive-field values.
Typical choices include:

\begin{itemize}
    \item a vector space $V$,
    \item a tensor space $T^r_s(X_{42})$,
    \item a Lie algebra $\mathfrak{g}$,
    \item a finite-dimensional manifold encoding interpretive states.
\end{itemize}

The specific structure of $F$ is not required for the algebraic
invariance of $I_p$, only that $F$ is smooth.

\subsection*{R5.3 Bundle Definition}

Define the interpretive-field bundle:
\[
\pi : E \to X_{42},
\]
where $E$ is a smooth manifold and $\pi$ is a smooth surjective map.

For each $x \in X_{42}$, the fiber is:
\[
E_x = \pi^{-1}(x) \cong F.
\]

\subsection*{R5.4 Local Triviality}

For each $x \in X_{42}$, there exists an open neighborhood $U_x$ and a
diffeomorphism:
\[
\Phi_x : \pi^{-1}(U_x) \to U_x \times F
\]
such that:
\[
\pi = \mathrm{proj}_{U_x} \circ \Phi_x.
\]

This ensures that $E$ is a smooth fiber bundle.

\subsection*{R5.5 Transition Functions}

For overlapping trivializations $(U_i, \Phi_i)$ and $(U_j, \Phi_j)$, the
transition map:
\[
t_{ij} : U_i \cap U_j \to \mathrm{Diff}(F)
\]
is defined by:
\[
\Phi_j \circ \Phi_i^{-1}(x, f) = (x, t_{ij}(x)(f)).
\]

Smoothness of $t_{ij}$ ensures that $E$ is a smooth bundle.

\subsection*{R5.6 Interpretive-Field Sections}

A smooth interpretive-field operator is a section:
\[
\mathcal{I}_{\Omega} \in \Gamma(E),
\]
where:
\[
\Gamma(E) = \{ s : X_{42} \to E \mid \pi \circ s = \mathrm{id}_{X_{42}} \}.
\]

Thus:
\[
\mathcal{I}_{\Omega}(x) \in E_x.
\]

\subsection*{R5.7 Interaction with Operators}

For any smooth operator $A : X_{42} \to X_{42}$, define the induced map:
\[
A^* : \Gamma(E) \to \Gamma(E)
\]
by:
\[
(A^* s)(x) = s(A(x)).
\]

This ensures compatibility between interpretive fields and the operator
algebra.

\subsection*{R5.8 Compatibility with the Interpretation Operator}

Since $I_p : X_{42} \to \{c\}$ is constant, it induces a trivial pullback
bundle:
\[
I_p^* E \cong \{c\} \times F.
\]

Thus:
\[
I_p(\mathcal{I}_{\Omega}(T(x))) = I_p(x)
\]
is well-defined because both sides lie in the same fiber over $c$.

\subsection*{R5.9 Tensor Interaction}

For any tensor operator $A : X_{42} \to X_{42}$, define:
\[
(I_p \otimes A)(x) = I_p(x).
\]

Since $I_p$ is constant, tensor interactions collapse to the fixed point.

\subsection*{R5.10 Summary}

The interpretive-field bundle $E$ is a smooth fiber bundle over the
42-dimensional interpretation manifold $X_{42}$, with:

\begin{itemize}
    \item smooth fibers $F$,
    \item smooth local trivializations,
    \item smooth transition functions,
    \item smooth sections representing interpretive fields,
    \item compatibility with all operators in the NDH algebra,
    \item a trivial pullback under the fixed-point operator $I_p$.
\end{itemize}

This structure ensures that interpretive-field operators are
mathematically admissible and that the fixed-point invariance of $I_p$
is well-defined.

This completes Appendix R5.
